\section{Sharp bounds for full displacement}
\label{sec:sharp-displacement}

Let $\sym{Usource},\sym{Vtarget}$ be subspaces of the same finite-dimensional real or complex
Hilbert space, with orthogonal projections $\sym{Psource},\sym{Qtarget}$. An admissible orthogonal
or unitary operator $\sym{Wcompetitor}$ satisfies $\sym{Wcompetitor}\sym{Usource}=\sym{Vtarget}$. We write $\sym{NormN}$ for a unitarily
invariant norm, or more generally a unitarily invariant seminorm. The
principal angles of $\sym{Usource}$ and $\sym{Vtarget}$ are $\like{theta}{\theta_j}\in[0,\pi/2]$; for an acute
pair let $\sym{alphaAngle}=\max_j\like{theta}{\theta_j}<\pi/2$, with $\sym{alphaAngle}=0$ when
$\sym{Usource}=\sym{Vtarget}=\{0\}$.

\begin{proposition}[Projection-distance lower bound]
\label{prop:projection-lower}
For every admissible $\sym{Wcompetitor}$ and every unitarily invariant seminorm $\sym{NormN}$,
\begin{equation}
  \sym{NormN}(\sym{Psource}-\sym{Qtarget})\leq \sym{NormN}(I-\sym{Wcompetitor}).
  \label{eq:projection-lower}
\end{equation}
\end{proposition}

\begin{proof}
Relative to $\sym{Usource}\oplus \like{Usource}{U^\perp}$, write
$\sym{Wcompetitor}=\left(\begin{smallmatrix}\sym{Ablock}&\sym{Bblock}\\\sym{Cblock}&\sym{Dblock}\end{smallmatrix}\right)$ and
$\sym{Jreflect}=2\sym{Psource}-I$. The off-diagonal part of $I-\sym{Wcompetitor}$ is
\[
 \frac{(I-\sym{Wcompetitor})-\sym{Jreflect}(I-\sym{Wcompetitor})\sym{Jreflect}}{2}
 =-\begin{pmatrix}0&\sym{Bblock}\\\sym{Cblock}&0\end{pmatrix}.
\]
By the triangle inequality and unitary invariance, its $\sym{NormN}$-seminorm is
at most $\sym{NormN}(I-\sym{Wcompetitor})$. Because $\sym{Wcompetitor}(\sym{Usource})=\sym{Vtarget}$, the singular values of $\sym{Ablock}$ are the
principal cosines $\cos\like{theta}{\theta_j}$. Unitarity gives
$\sym{Cblock}^*\sym{Cblock}=I-\sym{Ablock}^*\sym{Ablock}$ and $\sym{Bblock}\sym{Bblock}^*=I-\sym{Ablock}\sym{Ablock}^*$, so the nonzero singular values of both
$\sym{Bblock}$ and $\sym{Cblock}$ are the principal sines. The off-diagonal matrix therefore
has singular values $\sin\like{theta}{\theta_j}$ twice, with zeros for the remaining
dimensions. The difference $\sym{Psource}-\sym{Qtarget}$ has the same singular-value list
\citep{zhang2007angular,bhatia1997matrix}. Unitary invariance proves
\cref{eq:projection-lower}.
\end{proof}

\begin{theorem}[Sharp full-displacement comparison]
\label{thm:sharp-displacement}
Let $\sym{Usource},\sym{Vtarget}$ be an acute pair with largest principal angle $\sym{alphaAngle}$ and
canonical direct rotation $\sym{Rrotation}$. For every unitarily invariant seminorm $\sym{NormN}$
and every admissible $\sym{Wcompetitor}$,
\begin{equation}
 \boxed{\quad \sym{NormN}(I-\sym{Rrotation})\leq\sec(\sym{alphaAngle}/2)\,\sym{NormN}(\sym{Psource}-\sym{Qtarget})
                \leq\sec(\sym{alphaAngle}/2)\,\sym{NormN}(I-\sym{Wcompetitor}).\quad}
 \label{eq:sharp-displacement}
\end{equation}
For every fixed $0<\sym{alphaAngle}<\pi/2$, the coefficient $\sec(\sym{alphaAngle}/2)$
is best possible among coefficients valid for all dimensions, scalar
fields, and unitarily invariant norms. Consequently, the best common
coefficient under the hypothesis $\sym{alphaAngle}\leq\pi/3$ is $2/\sqrt3$;
$\sqrt2$ is a uniform bound for all acute pairs.
\end{theorem}

\begin{lstlisting}[style=leanpaper,
 caption={Lean signature for the sharp displacement bound, over real and
 complex finite-dimensional spaces for every unitarily invariant seminorm.
 The scalar-field and ambient-space context is omitted, and the linear
 isometry equivalence type is written without its Unicode notation.},
 label={lst:sharp-displacement}]
theorem cos_half_arcsin_projectionGap_mul_uiNorm_directRotation_le_displacement
    (N : UnitarilyInvariantSeminorm K E E)
    (U V : Submodule K E)
    [U.HasOrthogonalProjection] [V.HasOrthogonalProjection]
    (hacute : IsAcute U V)
    (W : LinearIsometryEquiv K E E)
    (hmap : U.map W.toLinearMap = V) :
    Real.cos (Real.arcsin (U.projectionGap V) / 2) *
      N (LinearMap.id - (directRotation U V hacute).toLinearMap) <=
        N (LinearMap.id - W.toLinearMap)
\end{lstlisting}

\begin{proof}
On a principal plane with angle $\like{theta}{\theta_j}$, the singular values of
$I-\sym{Rrotation}$ are $2\sin(\like{theta}{\theta_j}/2)$, twice, whereas those of $\sym{Psource}-\sym{Qtarget}$ are
$\sin\like{theta}{\theta_j}$, twice. Both operators vanish on the remaining fixed
summand. As $0\leq\like{theta}{\theta_j}\leq\sym{alphaAngle}<\pi/2$,
\[
 2\sin(\like{theta}{\theta_j}/2)
 =\frac{\sin\like{theta}{\theta_j}}{\cos(\like{theta}{\theta_j}/2)}
 \leq\sec(\sym{alphaAngle}/2)\sin\like{theta}{\theta_j}.
\]
Monotonicity of symmetric gauge functions gives the first inequality in
\cref{eq:sharp-displacement}; the second is
Proposition~\ref{prop:projection-lower}. For the family in
\cref{thm:small-angle}, both principal angles equal $\sym{alphaAngle}$ and
\[
 \|\sym{Psource}-\sym{Qtarget}\|_1=4\sin\sym{alphaAngle}
 =\|I-\like{Wcompetitor}{W_\alpha}\|_1,
 \qquad \|I-\like{Rrotation}{R_\alpha}\|_1=8\sin(\sym{alphaAngle}/2).
\]
Thus the ratio to the best admissible trace-norm displacement is exactly
$\sec(\sym{alphaAngle}/2)$. At $\sym{alphaAngle}=\pi/3$ it is $2/\sqrt3$, and it tends
to $\sqrt2$ as $\sym{alphaAngle}\uparrow\pi/2$.
\end{proof}

\begin{corollary}[Exact trace minimum for the real family]
\label{cor:real-trace-optimum}
For the subspaces $\sym{Usource},\like{Vtarget}{V_\theta}\subset\mathbb R^4$ of
\cref{thm:small-angle}, $0<\sym{theta}<\pi/2$,
\begin{equation}
 \min_{\sym{Wcompetitor}\in O(4):\,\sym{Wcompetitor}\sym{Usource}=\like{Vtarget}{V_\theta}}\|I-\sym{Wcompetitor}\|_1
 =\|\sym{Psource}-\sym{Qtarget}\|_1=4\sin\sym{theta},
 \label{eq:real-trace-minimum}
\end{equation}
with the minimum attained by the explicit $\like{Wcompetitor}{W_\theta}$ in
\cref{eq:wtheta}.
\end{corollary}

\begin{proof}
The singular-value calculation in \cref{thm:small-angle} gives
$\|\sym{Psource}-\sym{Qtarget}\|_1=4\sin\sym{theta}=\|I-\like{Wcompetitor}{W_\theta}\|_1$.
Apply Proposition~\ref{prop:projection-lower} to any other admissible orthogonal
map.
\end{proof}

The factor is optimal uniformly over unitarily invariant norms, not for
an individual norm: $\sym{NormQ}$-norms, for example, admit the stronger constant one.
The projection bound, comparison theorem, fixed-angle sharpness, and real
minimum are formalized in Lean; see Appendix~\ref{app:small-angle-artifact}.

The intertwining polar unitary itself is classical. Kato
\citep[Chapter~I, \S4.6, equations~(4.33)--(4.42)]{kato1966perturbation}
constructed it from the two projections, with earlier work by
Sz.-Nagy cited there. In the notation above the formula is
\[
 \sym{Rrotation}=\bigl(QP+(I-\sym{Qtarget})(I-\sym{Psource})\bigr)\bigl(I-(\sym{Psource}-\sym{Qtarget})^2\bigr)^{-1/2}.
\]
Murphy \citep[Lemma~6.2.1, p.~178]{murphy1990operator} established an
operator-norm comparison from the same polar factor. The estimate
\cref{eq:sharp-displacement} uses the projection singular values and
applies to every unitarily invariant seminorm; its global sharpness
comes from the real four-dimensional family. Bolotnikov, Li, and Rodman
\citep[Problem~1.1]{bolotnikov2001minimal} had already formulated the
corresponding ambient unitary-constrained norm optimization problem.

\subsection{Complex principal-plane minimizers}
\label{sec:complex-trace-model}

Davis and Kahan's complex phase construction
\citep[Example~4.2]{davis1970rotation} gives an exact trace minimum on
a principal plane. In suitable orthonormal coordinates, put
\begin{equation}
 \like{Rrotation}{R_\theta}=\begin{pmatrix}\cos\sym{theta}&i\sin\sym{theta}\\
                    i\sin\sym{theta}&\cos\sym{theta}\end{pmatrix},
 \qquad \like{Wcompetitor}{W_\theta}=e^{i\sym{theta}}\like{Rrotation}{R_\theta}.
 \label{eq:complex-phase}
\end{equation}
Both maps carry the first coordinate line to the same target line, while
$\like{Wcompetitor}{W_\theta}$ has eigenvalues $1,e^{2i\sym{theta}}$ and hence
$\|I-\like{Wcompetitor}{W_\theta}\|_1=2\sin\sym{theta}$ for $0\leq\sym{theta}\leq\pi/2$.
The projector difference on that plane has the same trace norm.

The formal development proves global trace-norm optimality for finite
orthogonal sums of these complex principal-plane models, including zero
and right angles, and with an unchanged orthogonal summand. The result is
preserved under an explicitly supplied ambient unitary equivalence.
For these models the minimum is $2\sum_j\sin\like{theta}{\theta_j}$. No acute-angle
assumption is needed. The distinction between this model theorem and an
all-pairs theorem is recorded in Appendix~\ref{app:small-angle-artifact}.

\ifincludeUnformalizedComplexTrace
\begin{tcolorbox}[
 enhanced,breakable,colback=white,colframe=red!80!black,
 colbacktitle=red!6!white,coltitle=red!80!black,
 title={Working extension --- not yet formalized for arbitrary complex pairs},
 fonttitle=\bfseries,boxrule=1pt,
 left=9pt,right=9pt,top=9pt,bottom=9pt,
 before skip=12pt,after skip=12pt]
\begin{proposition}[Complex trace minimum: proposed all-pairs formulation]
\label{prop:complex-trace-all-pairs}
For any equal-dimensional subspaces $\sym{Usource},\sym{Vtarget}$ of a finite-dimensional complex
Hilbert space, with principal angles $\like{theta}{\theta_1},\ldots,\like{theta}{\theta_k}$,
\[
 \min_{\sym{Wcompetitor}\text{ unitary}:\,\sym{Wcompetitor}\sym{Usource}=\sym{Vtarget}}\|I-\sym{Wcompetitor}\|_1
 =\|\sym{Psource}-\sym{Qtarget}\|_1=2\sum_{j=1}^k\sin\like{theta}{\theta_j}.
\]
\end{proposition}
\begin{proof}
The lower bound is Proposition~\ref{prop:projection-lower}. The cosine--sine
decomposition of $\sym{Usource},\sym{Vtarget}$ splits the nonzero principal angles into mutually
orthogonal complex two-planes, with a common fixed summand. On each plane
use the phase-adjusted map $\like{Wcompetitor}{W_{\theta_j}}$ from
\cref{eq:complex-phase}; it transports the corresponding source line
to its target line with trace displacement $2\sin\like{theta}{\theta_j}$. The identity
on the fixed summand contributes zero. The orthogonal sum is an
admissible unitary attaining $2\sum_j\sin\like{theta}{\theta_j}=\|\sym{Psource}-\sym{Qtarget}\|_1$.
\end{proof}

The Lean development proves the phase-plane construction and its
transport under a supplied coordinate equivalence. The general
principal-plane decomposition needed to instantiate that model for every
complex subspace pair remains to be formalized. This box is omitted from
the checked-only manuscript.
\end{tcolorbox}
\fi

\ifincludeUnformalizedRealTransport
\begin{tcolorbox}[
 enhanced,breakable,colback=white,colframe=red!80!black,
 colbacktitle=red!6!white,coltitle=red!80!black,
 title={Working extension --- real rank-one and codimension-one transport},
 fonttitle=\bfseries,boxrule=1pt,
 left=9pt,right=9pt,top=9pt,bottom=9pt,
 before skip=12pt,after skip=12pt]
\begin{proposition}[Exact real line/hyperplane norm minimization]
Let $n\geq2$ and let $\sym{Usource},\sym{Vtarget}\subset\mathbb R^n$ be lines or hyperplanes
at principal angle $\sym{theta}\in[0,\pi/2]$. Put
$a=2\sin(\sym{theta}/2)$ and, with trailing zero entries to size $n$, let
$\sym{Dplanar}=\operatorname{diag}(a,a,0,\ldots,0)$ and
$\sym{Hreflection}=\operatorname{diag}(2,0,\ldots,0)$. For every unitarily invariant
norm $\sym{NormN}$,
\[
  \min_{\sym{Wcompetitor}\in O(n):\,\sym{Wcompetitor}\sym{Usource}=\sym{Vtarget}} \sym{NormN}(I-\sym{Wcompetitor})
    =\min\{\sym{NormN}(\sym{Dplanar}),\sym{NormN}(\sym{Hreflection})\}.
\]
Both candidates are attained, respectively by a direct planar rotation
of determinant $+1$ and by a Householder reflection of determinant $-1$.
Consequently, Davis--Kahan's $\pi/3$ minimality claim holds for all
real subspace pairs in ambient dimension at most three. The verified
$\mathbb R^4$ counterexample shows that four is the minimum ambient
real dimension for failure of that short-angle assertion.
\end{proposition}
\begin{proof}
Select unit vectors $\sym{SourceVector},\sym{TargetVector}$ representing the source and target lines,
or unit normals for hyperplanes, so $\langle \sym{SourceVector},\sym{TargetVector}\rangle=\cos\sym{theta}$.
An admissible $\sym{Wcompetitor}$ has $\sym{Wcompetitor}\sym{SourceVector}=\pm \sym{TargetVector}$, whence
$\|I-\sym{Wcompetitor}\|_{\rm op}\geq 2\sin(\sym{theta}/2)=a$.
If $\det \sym{Wcompetitor}=1$, the nonreal eigenvalues of the real orthogonal matrix
occur in conjugate pairs and the multiplicity of $-1$ is even.
Thus the two largest singular values of $I-\sym{Wcompetitor}$ coincide, and both are
at least $a$. Fan dominance gives $\sym{NormN}(I-\sym{Wcompetitor})\geq \sym{NormN}(\sym{Dplanar})$, with equality
for the direct planar rotation. If $\det \sym{Wcompetitor}=-1$, $-1$ occurs as an
eigenvalue, and $I-\sym{Wcompetitor}$ has singular value $2$; hence
$\sym{NormN}(I-\sym{Wcompetitor})\geq \sym{NormN}(\sym{Hreflection})$. A single Householder reflection mapping $\sym{SourceVector}$
to $\sym{TargetVector}$ attains this bound (if $\sym{SourceVector}=\sym{TargetVector}$, choose a reflection fixing $\sym{SourceVector}$).
The minimum is therefore the smaller of the two values.
Finally, $\sym{Dplanar}$ is weakly majorized by $\sym{Hreflection}$ exactly when $2a\leq2$,
equivalently $\sym{theta}\leq\pi/3$. Every pair of nonzero proper
subspaces of the same rank in $\mathbb R^2$ or $\mathbb R^3$
is a line/hyperplane pair. The fixed-angle $\mathbb R^4$ family
from \cref{thm:small-angle} violates the asserted minimality at
arbitrarily small positive angles.
\end{proof}
The proof is conventional mathematics and has not been compiled as a
new Lean theorem. Its source-level novelty remains under investigation;
see the separate research note. This box is excluded from the
checked-only manuscript.
\end{tcolorbox}
\fi
